\subsection{Resultats}
\label{sec:DSCresultats}


Lors de la chauffe rapide (Fig. \ref{fig:chauffe_rapide}), nous constatons une température de transition vitreuse autour de 82 °C et un pic exothermique vers 244,8 °C d’une intensité de -43,15 J/g. En effet, chauffé à froid, les cristaux qui se défont libèrent une certaine chaleur. En revanche, le refroidissement rapide à chaud sur la Fig. \ref{fig:refroidissement_rapide} ne laisse pas au matériau le temps de se restructurer correctement et prend une structure amorphe.



\begin{flushleft}
\begin{minipage}{0.49\textwidth} 
\includegraphics[width=\textwidth]{dsc_chauffe.pdf}%
	\captionof{figure}{$\delta$H(T[$^\circ$C]) chauffe rapide}
       \label{fig:chauffe_rapide}
\end{minipage}
\hfill
\begin{minipage}{0.49\textwidth}
    \includegraphics[width=\linewidth]{dsc_amorphe.pdf}
	\captionof{figure}{$\delta$H(T[$^\circ$C]) refroidissement rapide}
        \label{fig:refroidissement_rapide}
\end{minipage}
\end{flushleft}
\\\\
Par la suite, le comportement de l’échantillon change lorsque nous lui permettons de se refroidir ou de chauffer plus lentement (visible sur les Figs. \ref{fig:chauffe_lente} et \ref{fig:refroidissement_lent}).

\begin{flushleft}
\begin{minipage}{0.49\textwidth}
    \includegraphics[width=\linewidth]{dsc_cristallisation_chaude.pdf}
        \captionof{figure}{$\delta$H(T[$^\circ$C]) chauffe lente}
        \label{fig:chauffe_lente}
\end{minipage}
\hfill
\begin{minipage}{0.49\textwidth}
    \includegraphics[width=\linewidth]{dsc_cristallisation_froide.pdf}
        \captionof{figure}{$\delta$H(T[$^\circ$C]) refroidissement lent}
        \label{fig:refroidissement_lent}
\end{minipage}
\end{flushleft}
Après un refroidissement rapide sur la Fig. \ref{fig:refroidissement_rapide}, nous chauffons lentement et obtenons la Fig. \ref{fig:chauffe_lente} sur laquelle $T_v$ se situe autour de 77 °C, soit environ 5 °C plus bas que lors de la première chauffe sur la Fig. \ref{fig:chauffe_rapide}. Nous observons ensuite un pic endothermique sur la Fig. \ref{fig:chauffe_lente} autour de 131 °C que nous ne voyons pas sur la Fig. \ref{fig:chauffe_rapide}. Cela démontre l’énergie de recristallisation du matériau chauffé à froid suite au refroidissement rapide ayant laissé une structure amorphe. Autour de 245 °C, avec 0,2 °C de différence par rapport à la chauffe rapide (Fig. \ref{fig:chauffe_rapide}), nous observons un second pic sur la Fig. \ref{fig:chauffe_lente}, qui est également de nature exothermique. Cela démontre l’énergie libérée lors de la décristallisation dans l’échantillon de PET. Lors du refroidissement lent sur la Fig. \ref{fig}, nous observons, contrairement à la Fig. \ref{fig}, un pic endothermique démontrant l’énergie nécessaire à la recristallisation du matériau


